\documentclass[10pt,a4paper]{amsart}
\usepackage[margin=19mm]{geometry}
\usepackage{amsmath,amssymb,mathtools}
\usepackage[scr=rsfs,cal=euler]{mathalfa}
\usepackage{xfrac}
\usepackage[unicode,hidelinks,bookmarks=false]{hyperref}
\newcommand{\ie}{i.e.\ }
\newcommand{\cf}{cf.\ }
\newcommand{\resp}{respectively\ }
\newcommand{\Subsection}{\subsection}
\providecommand{\nobreakdash}{\nobreak\hbox{-}}
\setlength{\emergencystretch}{2em}
\multlinegap0pt
\theoremstyle{plain}
\newtheorem{theorem}{Theorem}
\newtheorem{lemma}[theorem]{Lemma}
\newtheorem{proposition}[theorem]{Proposition}
\newtheorem{corollary}[theorem]{Corollary}
\theoremstyle{definition}
\newtheorem{definition}[theorem]{Definition}
\theoremstyle{remark}
\newtheorem{remark}[theorem]{Remark}
\usepackage{cleveref}
\makeatletter
\newcommand{\@currentthmtype}{}
\newcommand{\Pt}[1]{{\mathrm{Poly_{\partial}}({#1})}}
\newcommand{\Pext}[1]{{\mathrm{Poly}_{\mathrm{ext}}({#1})}}
\newcommand{\PPt}[2]{{\mathrm{PPoly}^{#1}_{\partial}({#2})}}
\newcommand{\Dt}[1]{{\mathrm{DivFree_{\partial}}({#1})}}
\newcommand{\mP}{\mathcal{P}}
\newcommand{\mQ}{\mathcal{Q}}
\newcommand{\R}{\mathbb{R}}
\makeatother
\title[Polynomial Interpolation of a Vector Field on a Convex Polyhedral Domain]{Polynomial Interpolation of a Vector Field on a Convex Polyhedral Domain: piecewise polynomial interpolation with tangency and divergence constraints (printed as Proposition 6.2)}
\author{Junyan Chu, Shizuo Kaji}
\date{Source: arXiv:2602.01803v2 (CC BY 4.0) (CC BY 4.0)}
\begin{document}
\maketitle

\noindent\emph{Source and licence.} Polynomial Interpolation of a Vector Field on a Convex Polyhedral Domain. Version arXiv:2602.01803v2 (CC BY 4.0), distributed under the Creative Commons Attribution 4.0 International licence, http://creativecommons.org/licenses/by/4.0/, as recorded in the arXiv licence metadata for this exact version and shown on the saved abstract page https://arxiv.org/abs/2602.01803v2. Full licence notice: licenses/arxiv-2602.01803-CC-BY-4.0.txt.\par\smallskip

\noindent\textit{Editorial scope.} Counted unit: the piecewise polynomial interpolation criterion with tangency and divergence constraints, printed in the article as Proposition 6.2 (source0.tex lines 776--778, one \texttt{proposition} environment with source label \texttt{prop:piecewise-interpolation}), delivered together with the article's complete original proof of it (source0.tex lines 779--793, one \texttt{proof} environment of 15 lines immediately adjacent to the statement, with no gap and no deferral). No mathematics is written, repaired or re-derived: statement and proof are the article's own text line for line, in the article's own notation, and no line of either is omitted, substituted or re-flowed.

Required context retained uncounted: the article's Proposition 6.1 on interpolation of prescribed vectors at prescribed points, together with its complete proof, because the delivered proof uses the Lagrange polynomials introduced there (source0.tex lines 649--665); the article's divergence-free interpolation proposition with the labelled lift formula that the delivered proof invokes (source0.tex lines 685--687 and 690--693); and the article's gluing proposition that supplies the matching condition (source0.tex lines 761--767). Those are exactly the labelled statements the delivered proof invokes, and nothing else from the article's development is used.

References. No citation command occurs in any retained line, so this package carries no bibliography block and no bibliography entry.

Editorial formatting only, in addition: an AMS \texttt{amsart} document that restates verbatim the article's own macro definitions needed by the retained lines, loads the standard cleveref package for the article's own cross-reference command, and restates the article's own theorem-environment declarations with the same shared counter the article declares. The article's own publisher class file \texttt{siamart251216.cls}, its bibliography style file and its bibliography database are neither shipped in this package nor used to build it. Consequently the retained environments print here in the order in which they occur, so the counted proposition prints as Proposition 4 whereas the article prints it as Proposition 6.2 under its per-section numbering; the three retained context statements print as Propositions 1, 2 and 3 under the same shared counter, the equations of the retained spans are numbered anew in order of appearance, and every cross-reference inside the retained text resolves to this wrapper's numbering. That is the only change to the numbering. No figure, no image-inclusion line and no drawing code occur in any retained line, so no artwork is delivered or needed, although the article contains figures elsewhere.

Not retained: the article's abstract, its introduction, its algebraic preliminaries, its examples and its bibliography; none of that material is used as a step of the delivered proof, and the delivered proof is complete as the author wrote it. The package ships the article's concatenated source verbatim as \texttt{sources/193809/source0.tex}, and every delivered excerpt is an exact line range of that file. No mathematical statement, hypothesis, estimate or conclusion has been rewritten, repaired or added, and printed slips, if any, are preserved literally.
\begin{proposition}\label{prop:interpolation_possible}
Let $x_1,\ldots,x_n\in\mP$ be distinct, with $n\ge1$. There exists $\xi\in\Pt{\mP}$ satisfying $\xi(x_s)=u_s$ for all $s$ if and only if $u_s\in T_{x_s}$. When this condition holds, one can take
\begin{equation}\label{eq:degree-bound}
\deg\xi\le K:=n-1+\max_s|J_s|,
\qquad J_s=\{i \mid h_i(x_s)\ne0\}.
\end{equation}
\end{proposition}

\begin{proof}
Necessity follows from the boundary condition. For sufficiency, choose $a\in\R^d$ such that $\lambda_s=\langle a,x_s\rangle$ are distinct. Set
\[
L_s(x)=\prod_{t\ne s}
\frac{\langle a,x\rangle-\lambda_t}{\lambda_s-\lambda_t},\qquad
W_s(x)=\prod_{i\in J_s}\frac{h_i(x)}{h_i(x_s)},\qquad
\xi(x)=\sum_{s=1}^n L_s(x)W_s(x)u_s.
\]
Since $L_s(x_t)=\delta_{st}$ and $W_s(x_s)=1$, the field interpolates the data. Each summand is tangent to every facet: it vanishes there if $i\in J_s$, and its constant direction $u_s$ is tangent otherwise. Its degree is at most $n-1+|J_s|$.
\end{proof}

\begin{proposition}\label{prop:divfree-interpolation}
Let $d\ge2$ and let $x_1,\ldots,x_n\in\mP$ be distinct, with $n\ge1$. If $u_s\in T_{x_s}$ for every $s$, there exists $\xi\in\Dt{\mP}_{\le 2K}$ satisfying $\xi(x_s)=u_s$, where $K$ is given by \eqref{eq:degree-bound}.
\end{proposition}

\begin{equation}\label{eq:divfree-lift}
\mathcal R_{s,u}(b)
=bu+\frac{(y_s\cdot\nabla b)u-(u\cdot\nabla b)y_s}{d-1}.
\end{equation}

\begin{proposition}\label{prop:gluedlinear}
Local fields $\xi_c\in\Pext{\mP_c}_{\le k}$ are the restrictions of a field in $\PPt{r}{\Delta}_{\le k}$ if and only if
\begin{equation}\label{eq:piecewise-matching}
h_{cc'}^{r+1}\mid\xi_c-\xi_{c'}
\qquad\text{componentwise for every shared facet.}
\end{equation}
\end{proposition}

\begin{proposition}\label{prop:piecewise-interpolation}
Fix $r\ge0$ and let $x_1,\ldots,x_n\in\mQ$ be distinct, with $n\ge1$. There exists a finite $k$ and a field $\xi\in\PPt{r}{\Delta}_{\le k}$ with $\xi(x_s)=u_s$ for every $s$ if and only if $u_s$ is tangent to every exterior facet containing $x_s$. If $d\ge2$, the field may additionally be chosen divergence-free on every cell.
\end{proposition}

\begin{proof}
Fix $s$ and let $\mathcal C_s=\{c\mid x_s\in\mP_c\}$. Let $\mathcal E_s$ be the exterior facets of these cells not containing $x_s$, and let $\mathcal F_s$ be the shared facets between a cell in $\mathcal C_s$ and a cell outside it. Their facet equations are nonzero at $x_s$, since each is a facet of a cell containing $x_s$ but does not itself contain $x_s$. Using the Lagrange polynomials $L_s$ from the proof of \Cref{prop:interpolation_possible}, set
\[
b_s(x)=L_s(x)\prod_{E\in\mathcal E_s}\frac{h_E(x)}{h_E(x_s)}
\prod_{F\in\mathcal F_s}\left(\frac{h_F(x)}{h_F(x_s)}\right)^{r+1},
\qquad
\eta_{s,c}=\begin{cases}
b_su_s,&c\in\mathcal C_s,\\
0,&c\notin\mathcal C_s.
\end{cases}
\]
The factors $h_E$ and the direction $u_s$ ensure exterior tangency, while the factors $h_F^{r+1}$ ensure matching by \Cref{prop:gluedlinear}. Since $\eta_s(x_t)=\delta_{st}u_s$, the sum $\sum_s\eta_s$ is the required interpolant.

For $d\ge2$, replace $b_su_s$ by $\mathcal R_{s,u_s}(b_s^2)$ from \eqref{eq:divfree-lift}. As in \Cref{prop:divfree-interpolation}, this preserves exterior tangency and the observation values and gives zero divergence. The components remain divisible by $h_F^{2r+1}$ on each $F\in\mathcal F_s$, since the lift differentiates $b_s^2$ at most once. As $2r+1\ge r+1$, matching is preserved as well.
\end{proof}

\end{document}
